
\section{胀开存在且射影}

\textbf{22.3.1.}\quad 现在是时候证明\trBlowup{}总是存在的了。我们将看到两种论证，它们各有各的启发。二者都将推出\trBlowup{}\trMorphism{}是\trProjective{}的，从而是\trQuasicompact、\trProper、\trFiniteType{}且\trSeparated{}的。特别地，若 $Y\to Z$ \trQuasicompact（相应地，\trProper、\trFiniteType、\trSeparated），则 $\mathrm{Bl}_X\,Y\to Z$ 亦复如是。（又若 $Y\to Z$ \trProjective{}而 $Z$ \trQuasicompact，则 $\mathrm{Bl}_X\,Y\to Z$ \trProjective。至于这条恼人的附加假设缘何出现，参看\mbox{练习~17.3.I} 的解答。）$k$-\trVariety{}的\trBlowup{}是 $k$-\trVariety（利用\trReducedness{}被保持这一事实，\mbox{练习~22.2.C}），而\trIrreducible{} $k$-\trVariety{}的\trBlowup{}是\trIrreducible{}的 $k$-\trVariety（利用\trIrreducibility{}被保持这一事实，同样见\mbox{练习~22.2.C}）。

\textbf{22.3.2. \trhTheorem（胀开的 $\mathit{Proj}$ 描述）。}\ ——\ \emph{设 $X\hookrightarrow Y$ 是由一个\trFiniteType\trQC\trIdealSheaf{} $\mathscr{I}\hookrightarrow\OO_Y$ 切出的\trClosedSubscheme。则}
\[
\mathit{Proj}_Y\bigl(\OO_Y\oplus\mathscr{I}\oplus\mathscr{I}^2\oplus\mathscr{I}^3\oplus\cdots\bigr)\longrightarrow Y
\]
\emph{满足胀开的\trUniversalProperty。}

（我们在\mbox{练习~15.6.E} 中赋予了\trIdealSheaf{}之乘积、从而 $\mathscr{I}^n$ 的意义。）
我们很快就会证明\mbox{定理~22.3.2}（\S22.3.3），不过先来看看它告诉了我们什么。由于 $\mathscr{I}$ 是有限型的，这个\trSheafOfGradedAlgebras{}的次数 $1$ 分片是有限型的。这个\trSheafOfGradedAlgebras{}显然也在次数 $1$ 处生成。于是这个\trSheafOfAlgebras{}满足\mbox{假设~17.2.1}（「在次数 $1$ 处有限生成」）。

但首先，我们应当确认 $X$ 的原像确实是一个\trEffectiveCartierDivisor。我们可以仿射局部地工作（\mbox{练习~22.2.A}），故不妨设 $Y=\Spec B$，而 $X$ 由有限生成的\trIdeal{} $I$ 切出。于是
\[
\mathrm{Bl}_X\,Y=\mathit{Proj}_B\bigl(B\oplus I\oplus I^2\oplus\cdots\bigr).
\]
我们稍稍滥用了记号：使用 $\mathrm{Bl}_X\,Y$ 这个记号，是因为我们尚未证明它满足\trUniversalProperty。回想（\S13.9.1），\trGradedRing{} $B\oplus I\oplus\cdots$——记作 $B(I)_\bullet$——称为 $B$ 中\trIdeal{} $I$ 的\trRees{}代数。

$X$ 的原像可不是随随便便一个\trEffectiveCartierDivisor{}就够了；它相应于这个 $\mathit{Proj}$ 上的\trInvertibleSheaf{} $\OO(1)$。的确，$\OO(1)$ 相应于如下操作：取我们的\trGradedRing，砍掉最底下的分片，再把全部分次片向左滑动 $1$ 位（\S15.1.4）；它就是相应于分次模
\[
IB(I)_\bullet=I\oplus I^2\oplus I^3\oplus\cdots
\]
（其中第一个直加项 $I$ 的分次为 $0$）的那个\trInvertibleSheaf。但它也可以被解读为 $X$ 的\trSchemeTheoreticPreimage，后者相应于 $B$ 的\trIdeal{} $I$：
\[
IB(I)_\bullet=I\bigl(B\oplus I\oplus I^2\oplus\cdots\bigr)\hookrightarrow B\oplus I\oplus I^2\oplus\cdots=B(I)_\bullet.
\]

于是，$X\hookrightarrow Y$ 在
\[
\mathit{Proj}\bigl(\OO_Y(\mathscr{I})_\bullet\bigr)=\mathit{Proj}\bigl(\OO_Y\oplus\mathscr{I}\oplus\mathscr{I}^2\oplus\cdots\bigr)
\]
中的\trSchemeTheoreticPreimage——它相应于 $\mathscr{I}\OO_Y(\mathscr{I})_\bullet=\mathscr{I}\oplus\mathscr{I}^2\oplus\mathscr{I}^3\oplus\cdots$ 所相应的那个\trIdealSheaf——是类 $\OO(1)$ 的一个\trEffectiveCartierDivisor。一旦我们验证了这套构造确实给出\trBlowup，这个除子就将是我们那个\trExceptionalDivisor{} $\mathrm{E}_X\,Y$。

此外，我们看到\trExceptionalDivisor{}可以被漂亮地描述为 $X$ 上的一个 $\mathit{Proj}$：
\[
\tag{22.3.2.1}\label{eq:exceptional-divisor-proj}
\mathrm{E}_X\,Y=\mathit{Proj}_X\bigl(\OO_Y/\mathscr{I}\oplus\mathscr{I}/\mathscr{I}^2\oplus\mathscr{I}^2/\mathscr{I}^3\oplus\cdots\bigr)=\mathit{Proj}_X\bigl(\OO_Y(\mathscr{I})_\bullet/(\mathscr{I}\OO_Y(\mathscr{I})_\bullet)\bigr).
\]
日后我们会看到，在好的境况下（若 $X$ 是 $Y$ 中的\trRegularEmbedding），这是某个\trVectorBundle{}的射影化（「\trProjectivization\trNormalBundle」）；见\mbox{练习~22.3.D}(a)。

\textbf{22.3.3. \mbox{定理~22.3.2} 的证明。}\quad 化归到仿射靶 $Y=\Spec B$（$X$ 相应于\trIdeal{} $I\subset B$）的情形。再化归到仿射源 $\Spec R$（带主\trEffectiveCartierDivisor{} $V(t)$，$t$ 是 $R$ 中的\trNonZerodivisor）的情形。（一个主\trEffectiveCartierDivisor{}局部地由单个\trNonZerodivisor{}切出。）于是我们化归到了 $\pi\colon\Spec R\to\Spec B$（相应于 $\phi\colon B\to R$）的情形。选取 $x_1,\dots,x_n$，使得 $(x_1,\dots,x_n)=I$，且 $(\phi(x_1),\dots,\phi(x_n))=(t)$。（这里我们用到了「$V(I)$ 的概形论拉回是 $V(t)$」这条假设，以及\mbox{练习~10.2.B} 对概形论拉回的理解。）从而
\[
\tag{22.3.3.1}\label{eq:phin-generated-unit}
(\phi(x_1)/t,\dots,\phi(x_n)/t)=R.
\]
（因为 $t$ 是\trNonZerodivisor，$\phi(x_i)/t$ 是有意义的。）我们将描述\emph{一个}映射 $\Spec R\to\Proj B(I)_\bullet$，它将把\trOpenSet{}上的那个映射 $\Spec R_t\to\Spec B$ 延拓开来。而由\mbox{练习~11.3.G}，它便是唯一的。考虑环映射 $\phi'\colon B(I)_\bullet\to R$，它把次数为 $d$ 的元素 $X$ 送到 $\phi(X)/t^d$。这就诱导出从 $\Spec R$ 到\textbf{\trAffine\trCone}（affine cone）$\Spec B(I)_\bullet$ 的一个映射 $\pi'$。

令 $X_i$ 为 $B(I)_\bullet$ 中相应于 $x_i$ 的次数 $1$ 元素。于是 $\phi'$ 诱导出
\[
\bigl((B(I)_\bullet)_{X_i}\bigr)_0\longrightarrow R_{\phi(x_i)/t},
\]
从而在几何上诱导出 $\mathrm{D}(\phi(x_i)/t)\to\mathrm{D}(X_i)$：从 $\Spec R$ 的一个\trDistinguishedOpen{}到\trAffine\trCone $\Spec B(I)_\bullet$ 的一个\trDistinguishedOpen。当 $x$ 跑遍集合 $\{x_1,\dots,x_n\}$ 时，诸 $\mathrm{D}(\phi(x_i)/t)$ 便跑遍 $\Spec R$ 的一组特异开覆盖（由 \eqref{eq:phin-generated-unit}）。于是，$\pi'\colon\Spec R\to\Spec B(I)_\bullet$ 的像避开\trIrrelevantIdeal{} $(X_1,\dots,X_n)\subset B(I)_\bullet$ 的\trVanishingSet，故映射 $\pi'$ 穿过 $\Spec B(I)_\bullet\setminus V(X_1,\dots,X_n)$，从而映入 $\Proj B(I)_\bullet$（一切都与到 $\Spec B$ 的态射交换）：
\[
\begin{tikzcd}[row sep=2.8em, column sep=3.4em]
\Spec B(I)_\bullet & \Spec B(I)_\bullet\setminus V(X_1,\dots,X_n) \ar[l,hook] \ar[r] \ar[d] & \Proj B(I)_\bullet \ar[dl] \\
\Spec R \ar[u,"\pi'"] \ar[ur,dotted] \ar[r,"\pi"] & \Spec B &
\end{tikzcd}
\]
我们便构造出了想要的态射 $\Spec R\to\Proj B(I)_\bullet$，它「提升」$\pi$。\hfill$\square$

以下是由这套\trBlowup{}构造派生的一些应用与观察。首先，我们可以验证，最初那批动机例子确实是胀开。例如，在原点处胀开 $\mathbb{A}^2$（带\trCoordinate{} $x$、$y$）给出：$B=k[x,y]$、$I=(x,y)$，而 $\Proj(B\oplus I\oplus I^2\oplus\cdots)=\Proj B[X,Y]$，其中 $B$ 的元素次数为 $0$，而 $X$ 与 $Y$ 的次数为 $1$，且分别与 $x$、$y$「相应」。

\textbf{22.3.4.}\ \emph{\trExceptionalDivisor{}的\trNormalBundle。}\quad 一个\trEffectiveCartierDivisor{} $D$ 的\trNormalBundle，就是（与）\trInvertibleSheaf{} $\OO(D)|_D$ 相伴随的空间——$\OO(D)$ 是\trTotalSpace{}上相应于 $D$ 的那个\trInvertibleSheaf，再限制到 $D$ 上（\mbox{练习~21.2.J}）。于是，在平面内一点处胀开的情形，\trExceptionalDivisor{}的\trNormalBundle{}是 $\OO(-1)$。（插一句：\trCastelnuovo{}判别准则（\mbox{定理~28.6.1}）断言其逆：给定一个光滑\trSurface，若它含有满足 $E\cong\mathbb{P}^1$ 且\trNormalBundle{}为 $\OO(-1)$ 的 $E$，则 $E$ 可以被收缩成另一个光滑\trSurface{}上的一点。）在正则\trVariety{}的正则\trSubvariety{}之胀开的情形，胀开结果是正则的（见\mbox{定理~22.3.10}），\trExceptionalDivisor{}是 $X$ 上的一个射影丛，而\trExceptionalDivisor{}的\trNormalBundle{}限制为 $\OO(-1)$（\mbox{练习~22.3.D}）。

\textbf{22.3.A.（较难而富启发的练习）}\quad 设 $\iota\colon X\hookrightarrow\mathbb{P}^n$ 是一个\trProjective\trScheme，请把 $X$ 上的\trAffine\trCone（\S9.3.11）在原点处的\trBlowup{}的\trExceptionalDivisor{}与 $X$ 本身等同起来，并证明它的\trNormalBundle（\S22.3.4）\trIsomorphic{}于 $\OO_X(-1):=\iota^*\OO_{\mathbb{P}^n}(-1)$。（在 $X=\mathbb{P}^1$ 的情形，我们便重新得到了平面在一点处的胀开。特别地，我们重新得到了那条重要事实：\trExceptionalDivisor{}的\trNormalBundle{}是 $\OO(-1)$。）

\textbf{22.3.5.}\ \emph{\trNormal\trCone。}\quad 受 \eqref{eq:exceptional-divisor-proj} 以及下方\mbox{练习~22.3.D} 的激励，我们给出如下定义。若 $X$ 是 $Y$ 中由 $\mathscr{I}$ 切出的\trClosedSubscheme，则 $X$ 在 $Y$ 中的\textbf{\trNormal\trCone}（normal cone）$\mathrm{N}_X\,Y$ 定义为
\[
\tag{22.3.5.1}\label{eq:normal-cone}
\mathrm{N}_X\,Y:=\mathit{Spec}_X\bigl(\OO_Y/\mathscr{I}\oplus\mathscr{I}/\mathscr{I}^2\oplus\mathscr{I}^2/\mathscr{I}^3\oplus\cdots\bigr).
\]
这可以卓有成效地解读为无穷小「\trTubularNeighborhood」的一种代数几何版本。但有几点当心事项。若 $Y$ \trSmooth，$\mathrm{N}_X\,Y$ 未必\trSmooth。（你可以自己算一算 $Y=\mathbb{A}^2_k$、$X=V(xy)$ 这个例子。）

而且即便 $X$ 与 $Y$ 都\trSmooth，虽然 $\mathrm{N}_X\,Y$ 是\trSmooth{}的（我们稍后就会看到，\mbox{练习~22.3.D}），它却无论如何也不「嵌入」到 $Y$ 之中。举一个简单的例子：考虑非有理\trCurve{} $Y=C$ 的一个\trSmooth{}点 $X=p$。那么 $\mathrm{N}_X\,Y$ 是有理的，从而不可能开嵌入到 $Y$ 中。

\textbf{22.3.6.}\ \emph{\trTangent\trCone。}\quad 若 $X$ 是一个\trClosedPoint{} $p$，则\trNormal\trCone{}称作 $Y$ 在 $p$ 处的\textbf{\trTangent\trCone}（tangent cone）。\textbf{\trProjectivization\trTangent\trCone}（projectivized tangent cone）则是\trExceptionalDivisor{} $\mathrm{E}_X\,Y$（同一个\trSheafOfGradedAlgebras{}的 $\mathit{Proj}$）。沿循 \S9.3.12，\trTangent\trCone{}与\trProjectivization\trTangent\trCone{}可以拼合到\trTangent\trCone{}的\trProjectiveCompletion{}之中：后者包含\trTangent\trCone{}作为\trOpenSubset，并包含\trProjectivization\trTangent\trCone{}作为互补的\trEffectiveCartierDivisor。

在\mbox{练习~22.3.D} 中我们将看到：在 $Y$ 的正则点处，\trTangent\trCone{}可以与\trTangentSpace{}等同起来，而\trNormal\trCone{}往往可以与\trNormalBundle{}的\trTotalSpace{}等同起来。

\textbf{22.3.B.（练习）}\quad 设 $Y=\Spec k[x,y]/(y^2-x^2-x^3)$（\trNode{}三次曲线；见图~\ref{fig:normalization-nodal-cubic} 的底部）。假设（以免分心）$\operatorname{char}k\neq 2$。证明：$Y$ 在原点处的\trTangent\trCone\trIsomorphic{}于 $\Spec k[x,y]/(y^2-x^2)$。于是，不严格地说，\trTangent\trCone「看起来就像」原本那个\trVariety{}被「无穷放大」了。

\textbf{22.3.C.（练习）}\quad 设 $S_\bullet$ 是域 $k$ 上一个\trFinitelyGenerated{}的\trGradedAlgebra。\mbox{练习~22.3.A} 给出了 $\Proj S_\bullet$ 与「$\Spec S_\bullet$ 在原点处的胀开」的\trExceptionalDivisor{}之间的一个\trIsomorphism。证明：$\Spec S_\bullet$ 在原点处的\trTangent\trCone\trIsomorphic{}于 $\Spec S_\bullet$ 本身。（你的几何直觉应当会使你觉得这些事实可信。）

\textbf{22.3.7. 沿正则嵌入胀开。}\quad 对正则嵌入 $X\subset Y$ 作胀开，这一情形特别令人愉悦。例如，\trExceptionalDivisor{}是 $X$ 上的一个射影丛。

\textbf{22.3.8. \trhTheorem。}\ ——\ \emph{若 $I\subset A$ 由一个\trRegularSequence{} $a_1,\dots,a_d$ 生成，则自然映射 $\Sym^n_A(I/I^2)\to I^n/I^{n+1}$ 是一个\trIsomorphism。}

\textbf{22.3.9. \trhCorollary。}\ ——\ \emph{若\trClosedEmbedding{} $\iota\colon X\hookrightarrow Y$ 是带\trIdealSheaf{} $\mathscr{I}\subset\OO_Y$ 的一个\trRegularEmbedding，则自然映射 $\Sym^n(\mathscr{I}/\mathscr{I}^2)\to\mathscr{I}^n/\mathscr{I}^{n+1}$ 是一个\trIsomorphism。此外，结合\mbox{命题~21.2.16}，我们看到 $\mathscr{I}^n/\mathscr{I}^{n+1}$ 是一个\trLocallyFreeSheaf。}

\mbox{练习~24.6.E} 在 $A$ 是 $k$-代数的情形给出了\mbox{定理~22.3.8}（以及\mbox{推论~22.3.9}）的一个稍显简单的证明。

在 \S22.3.11 开始\mbox{定理~22.3.8} 的证明之前，我们先展示它的用处。

\textbf{22.3.D.（练习）}\quad
（a）设 $X\to Y$ 是带\trIdealSheaf{} $\mathscr{I}$ 的一个\trRegularEmbedding。请把正规层（「正规丛」）的\trTotalSpace（\S17.1.5）与\trNormal\trCone $\mathrm{N}_X\,Y$（\eqref{eq:normal-cone}）等同起来，并证明\trExceptionalDivisor{} $\mathrm{E}_X\,Y$ 是 $X$ 上的一个射影丛（「\trProjectivization\trNormalBundle」）。

（b）证明：$\mathrm{E}_X\,Y$ 在 $\mathrm{Bl}_X\,Y$ 中的\trNormalBundle{}是 $\OO(-1)$（对这个 $X$ 上的射影丛而言）。

（c）再设 $X$ 是一个\trReduced{}的\trClosedPoint{} $p$。证明：$p$ 是 $Y$ 的一个正则点。请把 $p$ 的\trTangentSpace{}的\trTotalSpace{}与 $Y$ 在 $p$ 处的\trTangent\trCone{}等同起来。

\textbf{22.3.10. \trhTheorem。}\ ——\ \emph{设 $X\hookrightarrow Y$ 是 $k$ 上光滑\trVariety{}之间的一个\trClosedEmbedding。则 $\mathrm{Bl}_X\,Y$ 也是\trSmooth{}的。}

\emph{证明。}\quad 我们用到「光滑\trVariety{}是正则的」，此即\mbox{光滑性--正则性比较定理~13.2.7}(b)——其证明尚待完成。

我们不妨设 $k=\overline{k}$，因为把结论作\trBaseChange{}到\trAlgebraicClosure{}后依然成立。理由：\trSmoothness{}可以在\trBaseChange{}到 $\overline{k}$ 之后检验，\mbox{练习~13.2.H}。胀开构造与到\trAlgebraicClosure{}的\trBaseChange{}交换：你可以直接审视这个构造，或者等到能够援引\mbox{练习~24.1.P}(a) 中的推广。

由\mbox{观察~22.2.2}，只须检验 $\mathrm{Bl}_X\,Y$ 在 $\mathrm{E}_X\,Y$ 的诸点处的\trSmoothness。由\mbox{练习~13.2.M}(b)，$X\hookrightarrow Y$ 是一个\trRegularEmbedding。那么由\mbox{练习~22.3.D}(a)，$\mathrm{E}_X\,Y$ 是 $X$ 上的一个射影丛，从而是\trSmooth{}的，进而在其\trClosedPoint{}处正则。但 $\mathrm{E}_X\,Y$ 是 $\mathrm{Bl}_X\,Y$ 上的一个\trEffectiveCartierDivisor。于是由正则性的\trSlicing{}判别法（\mbox{练习~13.2.C}），$\mathrm{Bl}_X\,Y$ 在 $\mathrm{E}_X\,Y$ 的诸\trClosedPoint{}处正则，从而在 $\mathrm{E}_X\,Y$ 的一切点处\trSmooth。\hfill$\square$

\textbf{22.3.11.}\ \hard\ \emph{\mbox{定理~22.3.8} 的证明。}\quad 头一遍阅读时，把\mbox{定理~22.3.8} 的证明跳过去也完全合理。我们对\mbox{定理~22.3.8} 的证明沿循 [F2, A.6.1]，而后者又沿循 [Dav]。证明将在 \S22.3.13 完成。首先，令 $\alpha$ 为\trGradedRing{}之间的映射
\[
\alpha\colon (A/I)[X_1,\dots,X_d]\longrightarrow\oplus_{n=0}^{\infty}I^n/I^{n+1},
\]
它把 $X_i$ 送到 $a_i$ 在 $I/I^2$ 中的像。显然，$\alpha$ 是\trSurjectiveMap。

\textbf{22.3.E.（练习）}\quad 证明：\mbox{定理~22.3.8} 将由「$\alpha$ 是一个\trIsomorphism」这条陈述推出。

由于 $a_1$ 是\trNonZerodivisor，我们可以把 $A[a_2/a_1,\dots,a_d/a_1]$ 解读为\trTotalFractionRing（定义于 \S6.6.35）的一个子环。特别地，由于 $A$ 是其全分式环的子环，映射 $A\to A[a_2/a_1,\dots,a_d/a_1]$ 是\trInjectiveMap。定义
\[
\beta\colon A[T_2,\dots,T_d]\longrightarrow A[a_2/a_1,\dots,a_d/a_1]
\]
为 $T_i\mapsto a_i/a_1$。显然，映射 $\beta$ 是\trSurjectiveMap，而 $L_i:=a_1T_i-a_i$ 落在 $\ker(\beta)$ 之中。

\textbf{22.3.12. \trhLemma。}\ ——\ \emph{$\beta$ 的\trKernel{}是 $(L_2,\dots,L_d)$。}

\emph{证明。}\quad 我们对 $d$ 作归纳来证明这条结果。先考虑基底情形 $d=2$。设 $F(T_2)\in\ker\beta$，于是 $F(a_2/a_1)=0$。然后应用\trRemainderTheorem（Remainder Theorem）的算法，把 $a_1^{\deg F}F(T_2)$ 除以 $a_1T_2-a_2=L_2$，得
\[
\tag{22.3.12.1}\label{eq:division-remainder}
a_1^{\deg F}F(T_2)=G(T_2)(a_1T_2-a_2)+R,
\]
其中 $G(T_2)\in A[T_2]$，而 $R\in A$ 是余项。代入 $T_2=a_2/a_1$（并用「$A\to A[a_2/a_1]$ 为\trInjectiveMap」这一事实），便得 $R=0$。于是 $(a_1T_2-a_2)G(T_2)\equiv 0\pmod{(a_1^{\deg F})}$。利用「$a_2$ 模 $(a_1^{\deg F})$ 是\trNonZerodivisor」这一事实（因为 $a_1^{\deg F},a_2$ 是一个\trRegularSequence，\mbox{练习~9.5.E}），一次短归纳便表明 $G(T_2)$ 的系数必全被 $a_1^{\deg F}$ 整除。于是 $F(T_2)$ 被 $a_1T_2-a_2=L_2$ 整除，故 $d=2$ 的情形得证。

现在考虑一般情形 $d>2$，假定结果对一切更小的 $d$ 成立。令 $A'=A[a_2/a_1]$。那么 $a_1,a_3,a_4,\dots,a_d$ 是 $A'$ 中的一个\trRegularSequence。理由：$a_1$ 是 $A'$ 中的\trNonZerodivisor。$A'/(a_1)=(A[T_2]/(a_1T_2-a_2))/(a_1)=A[T_2]/(a_1T_2-a_2,a_1)=A[T_2]/(a_1,a_2)$。于是 $a_3$ 是这个环中的\trNonZerodivisor，因为它是 $A/(a_1,a_2)$ 中的\trNonZerodivisor。$a_4$ 是 $A[T_2]/(a_1,a_2,a_3)$ 中的\trNonZerodivisor，因为它是 $A'/(a_1,a_2,a_3)$ 中的\trNonZerodivisor，依此类推。正则序列之\mbox{定义~9.5.4} 的条件 (ii) 也成立，因为
\[
A'/(a_1,a_3,a_4,\dots,a_d)=A[T_2]/(a_1,a_2,\dots,a_d)\neq 0.
\]
考虑复合映射
\[
A[T_2,\dots,T_d]\to A'[T_3,\dots,T_d]\longrightarrow A'[a_3/a_1,\dots,a_d/a_1]=A[a_2/a_1,\dots,a_d/a_1].
\]
由 $d=2$ 的情形，第一个映射的\trKernel{}是 $L_2$。由归纳假设，第二个映射的\trKernel{}是 $(L_3,\dots,L_d)$。结果得证。\hfill$\square$

\textbf{22.3.13.}\ \emph{\mbox{定理~22.3.8} 的证明。}\quad 由\mbox{练习~22.3.E}，只须证明\trSurjectiveMap{} $\alpha$ 是一个\trIsomorphism。设 $F\in\ker(\alpha)$；我们想证明 $F=0$。不妨设 $F$ 是\trHomogeneous{}的，比如说次数为 $n$。考虑提升 $\alpha$ 的映射 $\alpha'\colon A[X_1,\dots,X_d]\to\oplus_{n=0}^{\infty}I^n/I^{n+1}$。把 $F$ 提升到 $A[X_1,\dots,X_d]$，于是 $F\in\ker(\alpha')$。我们想证明 $F\in IA[X_1,\dots,X_d]$。设 $F(a_1,\dots,a_d)=x\in I^{n+1}$。那么我们可以把 $x$ 写成 $F'(a_1,\dots,a_d)$，其中 $F'$ 是一个与 $F$ 同次数的\trHomogeneous{}多项式，其系数在 $I$ 中。于是把 $F$ 换成 $F-F'$，我们就被化归到下述问题：设 $F\in A[X_1,\dots,X_d]$ 是次数为 $n$ 的\trHomogeneous{}多项式，且 $F(a_1,\dots,a_d)=0$，我们想证明 $F\in IA[X_1,\dots,X_d]$\footnote{原书此处作「$IA[X_1,\dots,X_n]$」，当为「$IA[X_1,\dots,X_d]$」之误，径改。——译者注}。但若 $F(a_1,\dots,a_d)=0$，则 $F(1,a_2/a_1,\dots,a_d/a_1)=0$ 于 $A[a_2/a_1,\dots,a_d/a_1]$ 中。而\mbox{引理~22.3.12} 恰恰识别了 $\beta$ 的\trKernel，故
\[
F(1,T_2,T_3,\dots,T_d)\in(a_1T_2-a_2,\,a_1T_3-a_3,\,\dots,\,a_1T_d-a_d).
\]
于是 $F$ 的系数都在 $(a_1,\dots,a_d)=I$ 之中，正如所愿。\hfill$\square$
